\chapter{Лабораторная работа №5. Вычисление интеграла свертки}
\label{ch:lab5}

\section*{Задание}

В работе выполняются четыре раздела:
\begin{enumerate}
    \item Вычисление интеграла свертки двух временных функций аналитическим способом с заданием пределов интегрирования;
    \item Вычисление интеграла свертки двух прямоугольных импульсов (длительностью 1~с и 2~с) численным интегрированием, в том числе для каждого момента времени на интервале $(-1, 4)$ с шагом 0{,}05~с;
    \item Свертка прямоугольного импульса и гармонического колебания длительностью 4~с;
    \item Проверка связи свертки и преобразования Фурье: спектр свертки равен произведению спектров сигналов \cite{gonorovsky2006}.
\end{enumerate}

Программная реализация выполнена на языке Python (модули NumPy, Matplotlib), файл \texttt{lab5.py}. Интегрирование выполняется методом трапеций по сетке с шагом $\Delta\tau = 0{,}001$~с.

\section{Аналитическое вычисление интеграла свертки}

Заданы два прямоугольных импульса единичной амплитуды:
\begin{equation}
    s_1(t) =
    \begin{cases}
        1, & 0 \le t \le 1, \\
        0, & \text{иначе};
    \end{cases}
    \qquad
    s_2(t) =
    \begin{cases}
        1, & 0 \le t \le 2, \\
        0, & \text{иначе}.
    \end{cases}
    \label{eq:lab5:signals}
\end{equation}

Интеграл свертки сигналов определяется выражением \cite{gonorovsky2006}:
\begin{equation}
    y(t) = \int\limits_{-\infty}^{\infty} s_2(\tau)\, s_1(t - \tau)\, d\tau .
    \label{eq:lab5:conv}
\end{equation}

Функция $s_1(t - \tau)$ отлична от нуля при $t - 1 \le \tau \le t$, поэтому ненулевой участок подынтегрального произведения --- пересечение интервалов $[t-1, t]$ и $[0, 2]$:
\begin{equation}
    \max(0, t-1) \le \tau \le \min(2, t).
\end{equation}

Длина пересечения и есть значение интеграла (подынтегральная функция равна 1 на пересечении). По шагам получаем трапецеидальную свертку:
\begin{equation}
    y(t) =
    \begin{cases}
        0, & t < 0, \\
        t, & 0 \le t < 1 \quad \text{(нарастание)}, \\
        1, & 1 \le t \le 2 \quad \text{(плато, равно меньшей длительности)}, \\
        3 - t, & 2 < t \le 3 \quad \text{(спадание)}, \\
        0, & t > 3.
    \end{cases}
    \label{eq:lab5:trapezoid}
\end{equation}

Для типовых моментов $t$ пределы интегрирования и результат собраны в таблице:

\begin{table}[H]
    \caption{Аналитическое вычисление интеграла свертки по шагам}
    \centering
    \begin{tabular}{ccc}
        \toprule
        $t$, c & Пределы интегрирования $\tau$ & $y(t)$ \\
        \midrule
        $-0{,}5$ & пересечения нет & $0$ \\
        $0{,}5$  & $[0{,}00,\; 0{,}50]$ & $0{,}5$ \\
        $1{,}5$  & $[0{,}50,\; 1{,}50]$ & $1{,}0$ \\
        $2{,}5$  & $[1{,}50,\; 2{,}00]$ & $0{,}5$ \\
        $3{,}5$  & пересечения нет & $0$ \\
        \bottomrule
    \end{tabular}
\end{table}

\section{Численное интегрирование интеграла свертки}

Интегрирование \eqref{eq:lab5:conv} выполнено методом трапеций на сетке $\tau \in [-5, 5]$ c шагом 0{,}001~с. Замена переменной $t \to \tau$ и разворот функции $s_1(-\tau)$ показаны на рис.~\ref{fig:lab5:f1}.

\begin{figure}[H]
    \centering
    \includegraphics{lab5_fig1_s1_minus_tau.png}
    \caption{Исходные сигналы после замены переменной: $s_2(\tau)$ и развёрнутая функция $s_1(-\tau)$}
    \label{fig:lab5:f1}
\end{figure}

На рис.~\ref{fig:lab5:f2} показаны положения функции $s_1(t - \tau)$ для $t = -1, 0, 2$ и произведения $s_2(\tau)\, s_1(t-\tau)$; площадь каждого произведения равна значению интеграла свертки в данный момент $t$.

\begin{figure}[H]
    \centering
    \includegraphics{lab5_fig2_products.png}
    \caption{Сдвинутые функции $s_1(t-\tau)$, наложение на $s_2(\tau)$ и произведения для $t = -1, 0, 2$}
    \label{fig:lab5:f2}
\end{figure}

Значения интеграла для отдельных моментов времени (численное интегрирование):

\begin{table}[H]
    \caption{Результаты численного интегрирования для отдельных $t$}
    \centering
    \begin{tabular}{ccc}
        \toprule
        $t$, c & Аналитическое $y(t)$ & Численное $y(t)$ \\
        \midrule
        $0$    & $0$   & $0{,}001$ \\
        $0{,}2$ & $0{,}2$ & $0{,}200$ \\
        $1$    & $1$   & $1{,}001$ \\
        \bottomrule
    \end{tabular}
\end{table}

Расхождение в $\pm 0{,}001$ не превышает шага сетки и обусловлено методом трапеций на разрывах подынтегральной функции (границы импульсов).

Интеграл вычислен для каждого момента $t$ на интервале $(-1, 4)$ с шагом 0{,}05~с; результат вместе с аналитической трапецией \eqref{eq:lab5:trapezoid} изображен на рис.~\ref{fig:lab5:f3}. Максимальное отклонение численной свертки от аналитической составило $0{,}001$.

\begin{figure}[H]
    \centering
    \includegraphics{lab5_fig3_convolution_rects.png}
    \caption{Численная свертка прямоугольных импульсов (точки) и аналитический трапецеидальный импульс (штриховая линия)}
    \label{fig:lab5:f3}
\end{figure}

\section{Свертка прямоугольного сигнала и гармонического колебания}

В качестве $s_2(t)$ взято гармоническое колебание $s_2(t) = \cos(2\pi \cdot 0{,}5 \cdot t)$ длительностью 4~с (два периода), $s_1(t)$ --- прямоугольный импульс длительностью 1~с (рис.~\ref{fig:lab5:f4}, верхний график). Интеграл вычислен для каждого $t$ на интервале $(-1, 8)$ с шагом 0{,}05~с.

\begin{figure}[H]
    \centering
    \includegraphics{lab5_fig4_rect_cos.png}
    \caption{Исходные сигналы (вверху) и результат свертки импульса с гармоническим колебанием (внизу)}
    \label{fig:lab5:f4}
\end{figure}

Свертка с единичным окном есть скользящее среднее: при $1 \le t \le 4$
\begin{equation}
    y(t) = \int\limits_{t-1}^{t} \cos(\pi\tau)\, d\tau
         = \frac{\sin(\pi t) - \sin(\pi(t-1))}{\pi}
         = \frac{2}{\pi}\sin(\pi t),
\end{equation}
что соответствует виду графика: внутри «окна скольжения» амплитуда выходного колебания равна $2/\pi \approx 0{,}64$, на фронтах ($0 \le t \le 1$ и $4 \le t \le 5$) происходит плавный ввод и вывод сигнала.

\section{Связь свертки и преобразования Фурье}

Преобразование Фурье свертки двух сигналов равно произведению их преобразований Фурье: если $y(t) = s_1(t) * s_2(t)$, то
\begin{equation}
    Y(j\omega) = S_1(j\omega)\, S_2(j\omega), \qquad
    X(j\omega) = \int\limits_{-\infty}^{\infty} x(t)\, e^{-j\omega t}\, dt
    \label{eq:lab5:property}
\end{equation}
\cite{gonorovsky2006, oppenheim2012}.

Для сигналов из раздела 2 спектры $S_1(j\omega)$, $S_2(j\omega)$ и $Y(j\omega)$ вычислены численным интегрированием на интервале частот $-3 \le f \le 3$~Гц. Сравнение модуля и аргумента спектра свертки с произведением спектров приведено на рис.~\ref{fig:lab5:f5}.

\begin{figure}[H]
    \centering
    \includegraphics{lab5_fig5_ft_property.png}
    \caption{Проверка свойства \eqref{eq:lab5:property}: модуль (вверху) и аргумент (внизу) ПФ свертки и произведения $S_1(j\omega) \cdot S_2(j\omega)$}
    \label{fig:lab5:f5}
\end{figure}

Максимальная абсолютная погрешность совпадения кривых составила $0{,}0015$, относительная погрешность в зоне значимых значений спектра --- $0{,}07\,\%$. Свойство \eqref{eq:lab5:property} подтверждено с точностью численного интегрирования.

\endinput
